\chapter{Premiers pas}
\label{chap:premiers-pas}

\section{Premier lancement}

Au premier lancement, Pomone crée automatiquement une base de données SQLite vide à l'emplacement par défaut (voir chapitre \ref{chap:installation}). L'écran d'accueil affiche, une fois les premières plantations enregistrées, une \textbf{courbe d'occupation des planches} (surface cultivée plein champ et sous abri, semaine par semaine sur la saison) au-dessus d'une frise chronologique (Gantt) de la saison en cours. Voir la section \ref{chap:plantations} sur la \emph{vue Gantt} pour le détail.

\section{La barre latérale}

La navigation se fait via la barre latérale gauche qui reste visible en permanence :

\begin{itemize}
  \item \textbf{Accueil} : courbe d'occupation des planches et frise Gantt de la saison en cours.
  \item \textbf{Plantations} : la liste de vos plantations annuelles et pluriannuelles, avec frise Gantt détaillée.
  \item \textbf{Calendrier} : grille mensuelle unifiée mêlant tâches et jalons de culture (voir chapitre \ref{chap:calendrier}).
  \item \textbf{Tâches} : liste à plat de toutes les opérations, triée par date, avec badges \emph{En retard} / \emph{Aujourd'hui}.
  \item \textbf{Cultures} : les espèces et variétés que vous cultivez.
  \item \textbf{Lieux} : vos parcelles, planches, serres et autres emplacements.
  \item \textbf{Strates} : la hauteur de végétation (pour l'agroforesterie).
  \item \textbf{Paramètres} : configuration de la base de données et de la langue.
  \item \textbf{Carte} : vue d'occupation des lieux sur les 12 mois de la saison (voir chapitre \ref{chap:carte}).
\end{itemize}

\begin{notebox}
Chaque entrée de la barre latérale dispose d'un raccourci clavier \texttt{Ctrl+1} à \texttt{Ctrl+9} dans l'ordre indiqué ci-dessus, et \texttt{F1} ouvre ce manuel à tout moment.
\end{notebox}

\section{Découvrir Pomone avec un jeu de données de démo}

Pour explorer Pomone sans avoir à tout saisir manuellement, l'outil en ligne de commande \texttt{pomone-cli} peut peupler la base configurée avec un jeu de démonstration réaliste :

\begin{itemize}
  \item 5 cultures (Tomate, Carotte, Laitue, Courgette, Pommier)
  \item 7 variétés
  \item 7 lieux (un jardin avec 4 planches + 1 serre + un verger)
  \item 7 plantations (6 annuelles + 1 pluriannuelle) avec dates ancrées sur l'année courante
  \item 1 série récurrente (désherbage hebdomadaire)
\end{itemize}

Lancez la commande depuis la racine du dépôt (sur une base \textbf{vide}, sinon elle refuse de toucher à vos données) :

\begin{verbatim}
cargo run -p pomone-cli -- seed-demo
\end{verbatim}

La commande utilise la base configurée dans \texttt{\textasciitilde/.config/pomone/config.toml} (SQLite par défaut). Pour repartir sur une démo neuve, supprimez d'abord la base : \texttt{rm \textasciitilde/.local/share/pomone/pomone.sqlite}, puis relancez Pomone une fois pour recréer la base vide avant d'appeler \texttt{seed-demo}.

\section{Workflow recommandé pour démarrer}

\begin{enumerate}
  \item Vérifier les \textbf{strates} (chapitre \ref{chap:lieux}, section \ref{sec:strates}).
  \item Créer vos \textbf{lieux} : parcelles, puis planches à l'intérieur des parcelles (chapitre \ref{chap:lieux}).
  \item Ajouter vos \textbf{cultures et variétés} (chapitre \ref{chap:cultures}).
  \item Planifier vos \textbf{plantations} de la saison (chapitre \ref{chap:plantations}).
  \item Consulter le \textbf{calendrier} pour suivre les semis et récoltes (chapitre \ref{chap:calendrier}).
\end{enumerate}

% TODO capture — pré-prod : ajouter des captures d'écran annotées de la barre
% latérale et de l'écran d'accueil une fois en pré-production (données réelles).
